\documentclass[12pt,a4paper]{article}
\usepackage[utf8]{inputenc}
\usepackage[T1]{fontenc}
\usepackage[lithuanian]{babel}
\usepackage{lmodern}
\usepackage{amsmath, amssymb}
\usepackage{geometry}
\geometry{left=2cm, right=2cm, top=2.5cm, bottom=2.5cm}

\begin{document}

\begin{center}
    {\Large \textbf{Kontrolinis darbas: Daugianariai}}\\[0.3cm]
    {\large \textbf{Sokrato variantas -- Atsakymai ir sprendimai}}
\end{center}
\vspace{0.5cm}

\begin{enumerate}
    \item \textbf{Reiškinių tipai:}
    \begin{itemize}
        \item[a)] $\displaystyle \frac{x^3 - \sqrt{7}}{2}$ -- \textbf{R, S} (Racionalusis sveikas). Šaknis traukiama iš skaičiaus, kintamasis nėra vardiklyje.
        \item[b)] $\displaystyle \frac{3}{x^2 + 1}$ -- \textbf{R, T} (Racionalusis trupmeninis). Kintamasis yra vardiklyje.
        \item[c)] $\displaystyle x^{-3} + \pi x$ -- \textbf{R, T} (Racionalusis trupmeninis). Dėl neigiamo laipsnio rodiklio $x^{-3} = \frac{1}{x^3}$.
        \item[d)] $\displaystyle \sqrt{x^2 - 4}$ -- \textbf{I} (Iracionalusis). Kintamasis yra po šaknies ženklu.
        \item[e)] $\displaystyle 5x + 4^{-3}$ -- \textbf{R, S} (Racionalusis sveikas). Skaičius $4^{-3} = \frac{1}{64}$, kintamajam tai įtakos neturi.
        \item[f)] $\displaystyle \sqrt{3}x - 5$ -- \textbf{R, S} (Racionalusis sveikas). Šaknis traukiama tik iš skaičiaus $3$.
    \end{itemize}

    \vspace{0.3cm}
    \item \textbf{Vienanarių daugyba:} \\
    Sutvarkome reiškinį naudodami laipsnių savybes:
    $$ \left( -3x^3y^{-1} \right)^4 \cdot \left( 3x^{-2}y^2 \right)^{-3} = (-3)^4 \cdot x^{12} \cdot y^{-4} \cdot 3^{-3} \cdot x^6 \cdot y^{-6} $$
    $$ = 81 \cdot x^{12} \cdot y^{-4} \cdot \frac{1}{27} \cdot x^6 \cdot y^{-6} = 3 \cdot x^{12+6} \cdot y^{-4-6} = \mathbf{3x^{18}y^{-10}} \text{ arba } \mathbf{\frac{3x^{18}}{y^{10}}} $$
    \textbf{Paaiškinimas:} Tai \textbf{nėra vienanaris}, nes kintamojo $y$ laipsnio rodiklis yra neigiamas (tai reiškia dalybą iš kintamojo, kas prieštarauja vienanario apibrėžimui).

    \vspace{0.3cm}
    \item \textbf{Skaičiavimas be skaičiuotuvo:} \\
    \begin{align*}
        \frac{57^3+43^3}
        {0,1^{-2}\left(57\cdot14+43^2\right)}
        &= \frac{57^3+43^3}
        {100\left(57\cdot14+43^2\right)} \\
        &= \frac{(57+43)\left(57^2-57\cdot43+43^2\right)}
        {100\left(57\cdot14+43^2\right)} \\
        &= \frac{(57+43)\left(57(57-43)+43^2\right)}
        {100\left(57\cdot14+43^2\right)} \\
        &= \frac{100\left(57\cdot14+43^2\right)}
        {100\left(57\cdot14+43^2\right)} = 1.
    \end{align*}

    \vspace{0.3cm}
    \item \textbf{Formulių įrodymai:}
    \begin{itemize}
        \item[a)] $(x^r)^t = \underbrace{x^r \cdot x^r \cdot \dots \cdot x^r}_{t \text{ kartų}} = x^{\overbrace{r + r + \dots + r}^{t \text{ kartų}}} = \mathbf{x^{rt}}.$
        \item[b)] $(x+y)^3 = (x+y)(x+y)^2 = (x+y)(x^2 + 2xy + y^2) = x^3 + 2x^2y + xy^2 + x^2y + 2xy^2 + y^3 = \mathbf{x^3 + 3x^2y + 3xy^2 + y^3}.$
    \end{itemize}

    \vspace{0.3cm}
    \item \textbf{Didžiausia ir mažiausia reikšmės:} \\
    Išskiriame pilnąjį kvadratą:
    $$ -3x^2 + 2x - 1 = -3\left(x^2 - \frac{2}{3}x + \frac{1}{3}\right) = -3\left(x^2 - 2 \cdot x \cdot \frac{1}{3} + \frac{1}{9} - \frac{1}{9} + \frac{1}{3}\right) $$
    $$ = -3\left(\left(x - \frac{1}{3}\right)^2 + \frac{2}{9}\right) = -3\left(x - \frac{1}{3}\right)^2 - \frac{2}{3} $$
    Kadangi parabolės šakos nukreiptos žemyn ($a = -3 < 0$), reiškinys turi tik didžiausią reikšmę. \\
    \textbf{Didžiausia reikšmė:} $\mathbf{-\frac{2}{3}}$, ji gaunama, kai $x = \frac{1}{3}$. \\
    \textbf{Mažiausios reikšmės nėra} (artėja į $-\infty$).

    \vspace{0.3cm}
    \item \textbf{Sudėtingesnis reiškinio pertvarkymas:} \\
    Iškeliame bendrą dauginamąjį $(x-1)$:
    $$ (x-1)\left[ (x-1)^2 - (x-2)^2 \right] $$
    Išskleidžiame laužtiniuose skliaustuose esančius kvadratus (arba pritaikome kvadratų skirtumą):
    $$ (x-1) [ (x^2 - 2x + 1) - (x^2 - 4x + 4) ] = (x-1)[ 2x - 3 ] $$
    Atskliaudžiame: $2x^2 - 3x - 2x + 3 = 2x^2 - 5x + 3$. \\
    Išskiriame pilnąjį kvadratą:
    $$ 2\left(x^2 - \frac{5}{2}x\right) + 3 = 2\left(x^2 - 2 \cdot \frac{5}{4}x + \frac{25}{16} - \frac{25}{16}\right) + 3 = 2\left(x - \frac{5}{4}\right)^2 - \frac{25}{8} + \frac{24}{8} = 2\left(x - \frac{5}{4}\right)^2 - \frac{1}{8} $$
    Gavome $a = 2$, $b = -\frac{5}{4}$, $c = -\frac{1}{8}$. \\
    Skaičiuojame: $\sqrt{3(2) - 4(-\frac{5}{4}) - 8(-\frac{1}{8})+4} = \sqrt{6 + 5 + 1 + 4} = \sqrt{16} = \mathbf{4}$.

    \vspace{0.3cm}
    \item \textbf{Trijų kintamųjų mažiausia reikšmė:} \\
    Grupuojame narius ir išskiriame pilnuosius kvadratus:
    $$ (x^2 - 4x + 4) - 4 + (y^2 + 6y + 9) - 9 + (z^2 - 2z + 1) - 1 $$
    $$ = (x-2)^2 + (y+3)^2 + (z-1)^2 - 14 $$
    Kadangi bet kurio skaičiaus kvadratas yra $\ge 0$, \textbf{mažiausia reikšmė yra $\mathbf{-14}$} (ji įgyjama, kai $x=2, y=-3, z=1$).

\end{enumerate}

\end{document}